\section{Revisão da Literatura}

\subsection{Capital natural e condições produtivas}

A expansão agropecuária envolve uma escolha econômica entre usos alternativos da terra, na qual a conversão da vegetação pode ampliar a renda privada associada à produção, ao mesmo tempo em que reduz estoques de capital natural e serviços ecossistêmicos que nem sempre são incorporados aos retornos percebidos pelos agentes. Seus efeitos podem, portanto, ultrapassar os limites da propriedade convertida. No Cerrado brasileiro, medições em parcelas experimentais documentam diferenças de erosão e escoamento entre vegetação nativa, pastagem e cultivos \cite{Schwamback2024}. Na fronteira Amazônia--Cerrado, a substituição da floresta por soja foi associada a perdas de carbono até 100 cm do solo, com respostas distintas conforme a profundidade e o sistema produtivo; os estoques entre 300 e 800 cm mostraram-se relativamente pouco afetados pelo uso da terra \cite{Popin2025}.

Tais efeitos ajudam a exemplificar como decisões individualmente rentáveis podem modificar a base ambiental da produção. \citeonline{Brady2015} propõem uma abordagem de capital natural que relaciona serviços de suporte do solo e renda agrícola futura, utilizando funções de produção ajustadas a experimentos de longa duração. Nesse enquadramento, a contribuição ambiental interessa tanto ao rendimento quanto à eficiência dos demais insumos. A aproximação é pertinente ao problema deste artigo, sem implicar que carbono do solo e proporção de vegetação nativa sejam medidas intercambiáveis.

Com efeito, o problema permanece ainda menos resolvido quando se comparam benefícios e custos privados e públicos. Isso ocorre porque os serviços ambientais diferem entre territórios, operam em escalas distintas e muitos não possuem preços diretamente observáveis. No Pantanal, \citeonline{SeidlMoraes2000}, ao recalcularem para a Nhecolândia estimativas globais de serviços ecossistêmicos com dados locais, obtiveram valor por hectare aproximadamente metade daquele derivado dos parâmetros globais, diferença atribuída à heterogeneidade biofísica regional. Mais recentemente, \citeonline{FairbrassEtAl2025}, ao estruturarem indicadores de capital natural para o Pantanal, identificaram lacunas na disponibilidade e adequação dos dados necessários ao acompanhamento de seus ativos e serviços ambientais.

Esses trabalhos não identificam a interação entre vegetação e choque climático sobre o rendimento agrícola. Eles ajudam, contudo, a delimitar por que a contribuição do meio ambiente não deve ser presumida uniforme nem avaliada apenas por parâmetros externos ao território. Neste artigo, essa preocupação será traduzida em uma pergunta mais específica: se o estoque anterior de vegetação se relaciona à sensibilidade do rendimento aos choques de precipitação.

\subsection{Vegetação, clima e resposta da produtividade}

Não obstante, há já um acúmulo de evidências de que a vegetação compõe a base biofísica sobre a qual se desenvolve a atividade agrícola. \citeonline{SpracklenEtAl2012} relacionam a passagem de massas de ar sobre vegetação tropical à precipitação subsequente, evidenciando a relevância dos fluxos de umidade para além dos limites locais. \citeonline{JoshiGarkoti2023}, em outro contexto ecológico, documentam diferenças de propriedades e biomassa microbiana do solo entre tipos de vegetação no Himalaia central. A tradução desses mecanismos em resultados produtivos, todavia, mostra-se condicionada por características edáficas, pela composição e pelo arranjo espacial da paisagem e pelas práticas de manejo adotadas.

Tais elementos podem modificar tanto a sensibilidade da produção aos eventos climáticos quanto a estabilidade dos rendimentos. Para culturas nos Estados Unidos, \citeonline{HuangEtAl2021} mostram que textura, carbono orgânico, irrigação e profundidade da água subterrânea diferenciam as respostas ao clima. Na Inglaterra, \citeonline{RedheadEtAl2020} encontram relações entre habitats seminaturais e estabilidade ou resistência do rendimento do trigo. Já \citeonline{MorelEtAl2019}, ao examinar a cafeicultura de pequenos produtores no sudoeste da Etiópia durante o El Niño de 2015/2016, mostram que a vulnerabilidade se relaciona à interação entre características da paisagem e condições sociais. Registre-se que são culturas, escalas e desenhos empíricos distintos; a convergência sobre mecanismos não autoriza transportar seus coeficientes para Mato Grosso.

Porém, é muito importante ponderar que ainda não está estabelecido, no conjunto de estudos aqui examinado, em que medida as diferenças observadas na resposta produtiva aos eventos climáticos podem ser atribuídas ao estoque de vegetação previamente existente. A comparação permanece limitada porque parcela dessa literatura examina mecanismos biofísicos, associações espaciais ou sistemas produtivos particulares, empregando diferentes escalas, métricas ambientais e medidas de desempenho agrícola. Consequentemente, ainda é pouco claro quanto da heterogeneidade observada decorre da vegetação preexistente, em oposição a diferenças de solo, manejo, disponibilidade hídrica e demais características territoriais.

Essa distinção também se aplica às medidas de resistência ambiental. \citeonline{XiaoEtAl2023} identificam respostas distintas à seca e ao calor conforme a cobertura e o manejo, usando indicadores remotos de condição da vegetação. Esses indicadores permitem investigar resistência de ecossistemas, mas não medem diretamente o rendimento de soja em quilogramas por hectare. Sua contribuição para H2 é fundamentar a possibilidade de heterogeneidade biofísica, sem antecipar o sinal ou o tamanho do efeito agrícola que será estimado.

\subsection{Evidências brasileiras e heterogeneidade territorial}

No bioma Cerrado, sistemas de manejo que alteram a cobertura e o balanço hídrico do solo modificam o rendimento do milho e sua estabilidade sob condições climáticas distintas. \citeonline{SilvaEtAl2019} investigam esse mecanismo por meio de um modelo de crescimento de culturas calibrado e aplicado a condições contrastantes de clima e solo. Trata-se de evidência sobre agricultura conservacionista e cobertura do solo dentro do sistema de cultivo, e não de uma estimativa do efeito da proporção municipal de vegetação nativa.

Por exemplo, em Campo Verde, Mato Grosso, a comparação entre microbacias com Cerrado e pastagem identificou diferenças nas propriedades hidrofísicas do solo, na evapotranspiração e no escoamento \cite{NobregaEtAl2017}. Na Amazônia meridional mato-grossense, diferenças entre floresta e soja também foram observadas na precipitação que atravessa o dossel e alcança o solo \cite{BaseEtAl2012}. \citeonline{SchefflerEtAl2011}, por sua vez, identificam mudanças de infiltração e condutividade hidráulica entre floresta, pastagem e soja, mas concluem que, no ambiente estudado, o movimento predominantemente vertical da água permaneceu preservado. Este último resultado é especialmente útil por mostrar que alterações das propriedades do solo não implicam, necessariamente, a mesma resposta hidrológica em todos os contextos.

Como se vê, esses resultados conferem sustentação aos mecanismos pelos quais a condição ambiental preexistente pode modificar a forma como perturbações climáticas se transmitem à produção, mas não identificam, em um mesmo desenho empírico, quanto do efeito do choque sobre a produtividade agrícola depende especificamente do estoque anterior de vegetação nativa. No conjunto de evidências examinado, essa limitação também se manifesta no Pantanal. \citeonline{SeidlEtAl2001} investigam os incentivos econômicos associados ao desmatamento para implantação de pastagens cultivadas, identificando relações com riqueza fundiária e pecuária, esforço agrícola, população e acesso a infraestrutura e mercados. \citeonline{CardosoEtAl2011}, por sua vez, comparam vegetação arbórea nativa, pastagens cultivadas e regimes de pastejo no Pantanal sul-mato-grossense, documentando perda de matéria orgânica e redução da qualidade química do solo.

Os trabalhos avançam, portanto, sobre etapas distintas do problema --- determinantes da conversão e algumas de suas consequências biofísicas ---, mas não estabelecem como o estoque de vegetação \textit{anterior} a um choque climático modifica a resposta da produtividade a esse choque. Além disso, evidências obtidas em Mato Grosso do Sul não serão tomadas como estimativas representativas do Pantanal mato-grossense. A comparação demanda cuidado tanto com a localização quanto com a atividade econômica estudada: produtividade agrícola, efetivo bovino e qualidade de pastagens são desfechos diferentes.

De novo, Mato Grosso oferece, nesse sentido, um contexto especialmente informativo: a coexistência de Amazônia, Cerrado e Pantanal, associada a sistemas produtivos distintos, permite investigar, sob uma base territorial comum, se estoques previamente existentes de vegetação modificam a resposta da produtividade a choques climáticos comparáveis. A documentação da expansão e da intensificação agrícolas reforça a necessidade de considerar essas diferenças \cite{VanWeyEtAl2013,SimoesEtAl2020}. A viabilidade de uma comparação entre os três biomas, contudo, será demonstrada pelos dados, não deduzida da presença de seus limites no mapa estadual.

\subsection{Contribuição pretendida e limites da inferência}

Sob a perspectiva econômica, essa heterogeneidade importa porque alterações na sensibilidade da produtividade ao clima podem modificar, no nível do produtor, a variância dos rendimentos, o risco da atividade, a receita esperada, a rentabilidade e, consequentemente, decisões de investimento, manejo e uso da terra. O estoque de vegetação deixa, assim, de constituir apenas uma característica ambiental do território e passa a ser investigado como elemento potencialmente relevante para a formação dos resultados econômicos. As repercussões sobre lucro e decisões futuras constituem implicações a discutir; não serão confundidas com os desfechos diretamente observados neste artigo.

A literatura de clima e economia oferece uma base para examinar mudanças meteorológicas ao longo do tempo dentro da mesma unidade territorial \cite{DellEtAl2014}. Neste estudo, o interesse adicional recai sobre a diferença dessa resposta conforme o estoque ambiental anterior. Efeitos fixos ajudam a controlar características persistentes que afetam o nível do rendimento, mas não eliminam automaticamente diferenças não observadas na sensibilidade ao clima. Tampouco a defasagem da cobertura vegetal resolve sua seleção por produtores, tecnologias e condições locais.

Dito isso, a análise que se segue parte da caracterização dessas diferenças no território para, posteriormente, estimar as relações propostas, examinando não apenas efeitos médios, mas também magnitude econômica, incerteza e suporte das comparações entre biomas. Busca-se avançar sobre a ligação entre estado ambiental e resposta produtiva ao clima, mantendo aberta a possibilidade de resultados pequenos, imprecisos ou diferentes do mecanismo inicialmente considerado.
